\paragraph{El incremento de un menor doble.}
Fijemos un corte admisible \(b>0\) y \(j=\lfloor bn\rfloor\). En una máscara binaria regular
\(M\) de tamaño \(j\), escribamos \(A_R,A_C\) para las medidas uniformes sobre las filas y
columnas vivas. Sea \(\mathfrak F(M)\) el funcional de torres de la sección~\ref{sec:torres}.
Al derivarlo en sus medidas se mantiene la máscara de entrada y se recalcula su escalado
para que ambos márgenes sigan valiendo uno. Escribimos
\[
 h_R^M(i)=D_R\mathfrak F(M)[\delta_i-A_R],\qquad
 h_C^M(l)=D_C\mathfrak F(M)[\delta_l-A_C].
\]
Cada influencia tiene media cero respecto de su propia medida, por linealidad. Para el par
\(\parmar=\{e,f\}\) escribamos \(H_M(\parmar)=\sum_{z=e,f}(h_R^M(r_z)+h_C^M(c_z))\).
En el menor \(M_{\parmar}\) las medidas dan masa \(1/(j-2)\) a cada línea que queda. En particular,
\[
 A'_R-A_R=\frac{2A_R-\delta_{r_e}-\delta_{r_f}}{j-2},\qquad
 A'_C-A_C=\frac{2A_C-\delta_{c_e}-\delta_{c_f}}{j-2}.
\]
Ésta es una identidad exacta de medidas. Pongamos
\(\Delta_{\parmar}\mathfrak F=\mathfrak F(M_{\parmar})-\mathfrak F(M)\).
La cota uniforme de la segunda derivada en esas direcciones (Cauchy en las masas de las
líneas distinguidas) da, para todo par de la clase,
\[
 \Delta_{\parmar}\mathfrak F=-\frac{H_M(\parmar)}{j-2}+O(j^{-2}),\qquad
 j\Delta_{\parmar}\mathfrak F=-H_M(\parmar)+O(j^{-1}).
\]
Sean \(D\) el escalado bistocástico discreto de \(M\) y \(w_{\parmar}=(jD_e)(jD_f)\), acotado en la
clase. El cociente exacto de las dos probabilidades y su aproximación de torres satisfacen
\[
 \frac{\pi_{\parmar}^{\rm ex}}{\pi_{\parmar}^T}=\exp\{g_T(M_{\parmar})-g_T(M)\},\qquad
 \pi_{\parmar}^T=j^{-2}w_{\parmar}(1+O(j^{-1})).
\]
En efecto, la primera igualdad divide los cocientes de los permanentes y de sus predictores.
La segunda usa el término principal de la marginal de dos sitios y \(p_z/D_z=1+O(j^{-1})\).
Al restar \(\coefT(M_{\parmar})/(j-2)-\coefT(M)/j\), con \(\coefT=\mathfrak F-1/12\), los dos
incrementos sucesivos del resto de \((T^+)\) cuestan \(O(j^{-3})\). Así,
\[
 j^2\{g_T(M_{\parmar})-g_T(M)\}
 =j\Delta_{\parmar}\mathfrak F+2\mathfrak F(M)-1/6+O(j^{-1}).
\]
Desarrollar la exponencial del cociente exacto da entonces
\begin{equation}\label{eq:pt-par-fisico}
 -j^4(\pi_{\parmar}^{\rm ex}-\pi_{\parmar}^T)
 =w_{\parmar}\{1/6-2\mathfrak F(M)-j\Delta_{\parmar}\mathfrak F\}+O(j^{-1}).
\end{equation}
El resto es uniforme en los pares. Hay a lo sumo \(j^2(j-1)\) y la primera media de
\(\mathcal M_n(b)\) multiplica esta suma por \(n/j^4\). El coste total de ese resto y del
de sustituir \(-j\Delta_{\parmar}\mathfrak F\) por \(H_M(\parmar)\) es por tanto
\[
 O\!\left(\frac{n}{j^4}\,j^2(j-1)\,\frac1j\right)
 =O(n/j^2)=o_b(1).
\]
